\documentclass[12pt, letterpaper]{article}

% ── Packages ──────────────────────────────────────────────────────────────────
\usepackage[margin=1in]{geometry}
\usepackage[T1]{fontenc}
\usepackage[utf8]{inputenc}
\usepackage{lmodern}
\usepackage{microtype}
\usepackage{booktabs}
\usepackage{array}
\usepackage{longtable}
\usepackage{xcolor}
\usepackage{listings}
\usepackage{enumitem}
\usepackage{titlesec}
\usepackage{fancyhdr}
\usepackage[most,breakable]{tcolorbox}
\usepackage{hyperref}
\usepackage{parskip}
\usepackage{amsmath}
\usepackage{amssymb}
\setlength{\headheight}{14pt}

% ── Colors ────────────────────────────────────────────────────────────────────
\definecolor{warnorange}{RGB}{255,140,0}
\definecolor{dangerred}{RGB}{200,30,30}
\definecolor{noteblue}{RGB}{30,100,200}
\definecolor{codegray}{RGB}{245,245,245}
\definecolor{darkgray}{RGB}{60,60,60}

% ── Hyperref ──────────────────────────────────────────────────────────────────
\hypersetup{
    colorlinks=true,
    linkcolor=noteblue,
    urlcolor=noteblue,
    pdftitle={Wireless Drive-By-Wire User Manual},
    pdfauthor={Capstone Spring 2026}
}

% ── Code listings ─────────────────────────────────────────────────────────────
\lstdefinestyle{terminal}{
    backgroundcolor=\color{codegray},
    basicstyle=\ttfamily\small,
    breaklines=true,
    frame=single,
    rulecolor=\color{gray},
    framesep=6pt,
    xleftmargin=8pt,
    xrightmargin=8pt,
    aboveskip=8pt,
    belowskip=8pt
}
\lstdefinestyle{cpp}{
    backgroundcolor=\color{codegray},
    basicstyle=\ttfamily\small,
    keywordstyle=\color{blue}\bfseries,
    stringstyle=\color{orange!80!black},
    commentstyle=\color{green!50!black}\itshape,
    language=C++,
    breaklines=true,
    frame=single,
    rulecolor=\color{gray},
    framesep=6pt,
    xleftmargin=8pt,
    xrightmargin=8pt,
    aboveskip=8pt,
    belowskip=8pt
}

% ── Callout boxes ─────────────────────────────────────────────────────────────
\newtcolorbox{warningbox}{
    colback=warnorange!12,
    colframe=warnorange,
    fonttitle=\bfseries,
    title={$\bigtriangleup$\ WARNING},
    breakable
}
\newtcolorbox{dangerbox}{
    colback=dangerred!10,
    colframe=dangerred,
    fonttitle=\bfseries,
    title={!! DANGER},
    breakable
}
\newtcolorbox{notebox}{
    colback=noteblue!8,
    colframe=noteblue,
    fonttitle=\bfseries,
    title={Note},
    breakable
}

% ── Header / Footer ───────────────────────────────────────────────────────────
\pagestyle{fancy}
\fancyhf{}
\fancyhead[L]{\small\textcolor{darkgray}{Wireless Drive-By-Wire}}
\fancyhead[R]{\small\textcolor{darkgray}{User Manual}}
\fancyfoot[C]{\small\thepage}
\renewcommand{\headrulewidth}{0.4pt}

% ── Section formatting ────────────────────────────────────────────────────────
\titleformat{\section}{\large\bfseries}{}{0em}{\thesection\quad}[\vspace{2pt}\titlerule]
\titleformat{\subsection}{\normalsize\bfseries}{}{0em}{\thesubsection\quad}
\titleformat{\subsubsection}{\normalsize\itshape\bfseries}{}{0em}{\thesubsubsection\quad}
\setcounter{secnumdepth}{3}

% ══════════════════════════════════════════════════════════════════════════════
\begin{document}
% ══════════════════════════════════════════════════════════════════════════════

% ── Title Page ────────────────────────────────────────────────────────────────
\begin{titlepage}
  \centering
  \vspace*{2cm}

  {\Huge\bfseries Wireless Drive-By-Wire\\[0.4em]
  \Large User Manual}

  \vspace{0.5cm}\rule{\linewidth}{0.5pt}\vspace{1cm}

  {\large Capstone Spring 2026}\\[0.4cm]
  {\normalsize ECE 4773 --- Senior Design}

  \vspace{2cm}
  \begin{tabular}{ll}
    \textbf{Document Version:} & 1.0 \\
    \textbf{Date:}             & May 2026 \\
    \textbf{Status:}           & Final \\
  \end{tabular}

\end{titlepage}

\tableofcontents
\newpage

% ══════════════════════════════════════════════════════════════════════════════
\section{Introduction}
% ══════════════════════════════════════════════════════════════════════════════

The Wireless Drive-By-Wire (DBW) system allows an operator to remotely steer,
accelerate, and brake a Nissan Leaf using a Logitech G29 racing wheel and pedal set
connected to a laptop anywhere with an internet connection.

Control inputs are captured at \textbf{60\,Hz}, transmitted over UDP to a cloud relay
server hosted on Oracle Cloud Infrastructure (OCI), and forwarded to an ESP32
microcontroller installed on a custom PCB inside the vehicle.  The ESP32 runs a
closed-loop PID controller for steering---reading the actual steering angle from the
vehicle's CAN bus---and drives open-loop throttle and brake actuators through dual
MCP4728 digital-to-analog converters (DACs).

\subsection{System Architecture}

\begin{center}
\small\ttfamily
\begin{tabular}{ccccc}
  \textrm{[Driver Side]} &&\textrm{[Cloud]}&&\textrm{[Car Side]}\\[4pt]
  G29 + Laptop &$\xrightarrow{\text{UDP / internet}}$& OCI Relay Server
               &$\xrightarrow{\text{WiFi UDP}}$& ESP32 on PCB\\
  \small\textrm{g29\_to\_server\_original.py} && \small\textrm{relay\_server.py}
               && $\downarrow$\\
               &&&&\textrm{DAC outputs}\\
               &&&&$\downarrow$\\
               &&&&\textrm{Steer / Throttle / Brake}
\end{tabular}
\end{center}

The relay server is required because the ESP32 is behind a mobile hotspot with no
public IP.  The ESP32 registers itself via periodic \texttt{HEARTBEAT} packets;
the relay server records that public endpoint and uses it to forward incoming
\texttt{CMD} packets from the operator (UDP hole-punching).

% ══════════════════════════════════════════════════════════════════════════════
\section{Safety}
% ══════════════════════════════════════════════════════════════════════════════

\begin{dangerbox}
\begin{itemize}[leftmargin=*]
  \item A trained safety driver \textbf{must} be seated and belted in the vehicle at all times.
  \item The safety driver must be ready to take manual control \textbf{immediately}.
  \item \textbf{Never} operate on public roads.  Use only a closed, controlled area.
  \item \textbf{Never} exceed walking pace during initial testing.
  \item Keep all bystanders at least \textbf{50 feet} from the vehicle at all times.
\end{itemize}
\end{dangerbox}

\begin{warningbox}
\begin{itemize}[leftmargin=*]
  \item If the operator laptop stops sending (Ctrl+C or network loss), the ESP32
        has no watchdog --- it will continue targeting the \textbf{last commanded}
        steering angle, throttle, and brake indefinitely.  The safety driver must
        take manual control immediately if the connection is interrupted.
  \item Writing 0\,V to the brake DAC channels is \textbf{not} a neutral state and may
        be interpreted as a fault by the vehicle ECU.  The firmware always drives brake
        channels to their baseline voltages.
  \item Do not modify PID gains without re-validating on a stationary vehicle first.
\end{itemize}
\end{warningbox}

% ══════════════════════════════════════════════════════════════════════════════
\section{Hardware Overview}
% ══════════════════════════════════════════════════════════════════════════════

\subsection{Component List}

\begin{longtable}{>{\bfseries}p{4.6cm}p{1.2cm}p{7.4cm}}
  \toprule
  \textbf{Component} & \textbf{Qty} & \textbf{Notes}\\
  \midrule\endhead\bottomrule\endfoot
  Nissan Leaf (2015--2018)          & 1 & Vehicle under control; CAN ID \texttt{0x2} carries steering angle\\[3pt]
  Logitech G29 Wheel + Pedals       & 1 & USB to operator laptop\\[3pt]
  Operator Laptop                   & 1 & Any Python 3.8+ capable machine\\[3pt]
  ESP32 DevKit V1                   & 1 & Mounted on custom DBW PCB inside vehicle\\[3pt]
  Custom DBW PCB                    & 1 & Houses ESP32, CAN controller, and both DACs\\[3pt]
  MCP2515 CAN Controller            & 1 & On-board; reads steering angle via SPI\\[3pt]
  MCP4728 DAC $\times$2             & 2 & On-board; drives steering, throttle, brake\\[3pt]
  Mobile Hotspot                    & 1 & Provides WiFi to the ESP32 in the vehicle\\[3pt]
  OCI Relay Server                  & 1 & Oracle Cloud VM; UDP port 4210 must be open\\
\end{longtable}

\subsection{ESP32 Pin Reference}

\begin{center}
\begin{tabular}{ll}
  \toprule
  \textbf{Signal} & \textbf{ESP32 Pin}\\
  \midrule
  SPI SCK  & 18\\
  SPI MISO & 19\\
  SPI MOSI & 23\\
  CAN CS   & 5\\
  CAN INT  & 4\\
  I2C Bus~1 SDA / SCL (DAC~1) & 21 / 22\\
  I2C Bus~2 SDA / SCL (DAC~2) & 33 / 32\\
  \bottomrule
\end{tabular}
\end{center}

\subsection{DAC Channel Assignments}

\begin{center}
\begin{tabular}{llll}
  \toprule
  \textbf{DAC} & \textbf{Ch} & \textbf{Function} & \textbf{Baseline}\\
  \midrule
  DAC~2 & C & Steering main & 2.47\,V\\
  DAC~2 & D & Steering sub  & 2.47\,V\\
  DAC~1 & A & Throttle main & 0.37\,V\\
  DAC~1 & B & Throttle sub  & 0.75\,V\\
  DAC~1 & C & Brake main    & 3.42\,V\\
  DAC~1 & D & Brake sub     & 1.51\,V\\
  \bottomrule
\end{tabular}
\end{center}

\begin{notebox}
All DAC channels use the MCP4728 internal voltage reference
(\texttt{MCP4728\_VREF\_INTERNAL}) with $2\times$ gain, giving a full-scale output of
4.096\,V and 12-bit resolution ($\approx$1\,mV/step).
\end{notebox}

% ══════════════════════════════════════════════════════════════════════════════
\section{Software Installation}
% ══════════════════════════════════════════════════════════════════════════════

\subsection{Operator Laptop}

Install Python 3.8 or newer, then install the only required library:

\begin{lstlisting}[style=terminal]
pip install pygame
\end{lstlisting}

The server IP is already configured in \texttt{g29\_to\_server\_original.py}.

\subsection{OCI Relay Server}

The relay server uses only the Python standard library:

\begin{lstlisting}[style=terminal]
python3 relay_server.py
\end{lstlisting}

Verify that UDP port \textbf{4210} is open in both the OCI Security List and the
instance firewall (\texttt{iptables} / \texttt{firewalld}).

\subsection{ESP32 Firmware}

\subsubsection*{Required Tools}

\begin{itemize}
  \item Arduino IDE 2.x
  \item ESP32 board support package (install via \textbf{File\,$\to$\,Preferences\,$\to$\,Additional
        Board Manager URLs}, then search ``esp32'' in \textbf{Tools\,$\to$\,Board Manager})
  \item Library: \textbf{Adafruit MCP4728} (Tools\,$\to$\,Manage Libraries)
  \item Library: \textbf{mcp\_can} by Seeed Studio (Tools\,$\to$\,Manage Libraries)
\end{itemize}

\begin{notebox}
\texttt{WiFi} and \texttt{WiFiUdp} ship with the ESP32 board package---do \textbf{not}
install them separately.
\end{notebox}

\subsubsection*{Configure WiFi Credentials}

Open \texttt{esp32\_car.ino} and set the hotspot SSID and password:

\begin{lstlisting}[style=cpp]
const char* ssid     = "YOUR_HOTSPOT_NAME";
const char* password = "YOUR_HOTSPOT_PASSWORD";
\end{lstlisting}

\subsubsection*{Flash the Firmware}

\begin{enumerate}
  \item Connect the ESP32 to a computer via USB.
  \item \textbf{Tools\,$\to$\,Board}\,$\to$\,\textbf{ESP32 Dev Module}.
  \item \textbf{Tools\,$\to$\,Port}\,$\to$\,select the correct port
        (\texttt{COM3} on Windows, \texttt{/dev/ttyUSB0} on Linux).
  \item Click \textbf{Upload} and wait for ``Done uploading.''
\end{enumerate}

% ══════════════════════════════════════════════════════════════════════════════
\section{Pre-Operation Checklist}
% ══════════════════════════════════════════════════════════════════════════════

Complete every item in order before each operating session.

\begin{enumerate}[label=$\square$, leftmargin=2em]
  \item Vehicle is in a closed, controlled area with no bystanders within 50\,ft.
  \item Safety driver is seated, belted, and briefed.
  \item Mobile hotspot is powered on and broadcasting.
  \item Relay server is running on OCI (\lstinline[style=terminal]{python3 relay_server.py}).
  \item Relay server shows: \lstinline[style=terminal]{Relay server live on port 4210. Waiting for car heartbeat...}
  \item ESP32 is powered and WiFi-connected (Serial Monitor shows \lstinline[style=terminal]{WiFi connected}).
  \item Relay server shows: \texttt{[OK] CAR CONNECTED! Logged IP: ('x.x.x.x', 4210)}
  \item G29 is plugged in and has completed its self-calibration spin.
  \item \texttt{pygame} is installed on the operator laptop.
  \item Vehicle is in \textbf{Drive} with foot brake held by the safety driver.
\end{enumerate}

% ══════════════════════════════════════════════════════════════════════════════
\section{Operating the System}
% ══════════════════════════════════════════════════════════════════════════════

\subsection{Starting a Session}

Once all checklist items are satisfied, run the sender on the operator laptop:

\begin{lstlisting}[style=terminal]
python3 g29_to_server_original.py
\end{lstlisting}

Expected output:
\begin{lstlisting}[style=terminal]
Joystick detected: Logitech G29 Driving Force Racing Wheel
Sending UDP to OCI Server at 147.224.143.221:4210
seq=00000 | steer= 0.0000 | throttle=0.0000 | brake=0.0000
seq=00001 | steer= 0.0012 | throttle=0.0000 | brake=0.0000
...
\end{lstlisting}

The system is now live.

\subsection{Control Mapping}

\begin{center}
\begin{tabular}{lll}
  \toprule
  \textbf{G29 Input} & \textbf{Sent Value} & \textbf{Effect}\\
  \midrule
  Steering wheel left  & $-1.0$ to $0.0$ & Steer left (target up to $-250\,^\circ$)\\
  Steering wheel right & $0.0$ to $+1.0$ & Steer right (target up to $+250\,^\circ$)\\
  Throttle pedal       & $0.0$ to $1.0$  & Accelerate (open-loop DAC offset)\\
  Brake pedal          & $0.0$ to $1.0$  & Brake (open-loop DAC offset)\\
  \bottomrule
\end{tabular}
\end{center}

\begin{notebox}
The G29 pedal axes report \texttt{-1.0} when fully pressed and \texttt{+1.0} when
released.  The sender inverts this so that \texttt{1.0} always means fully pressed.
\end{notebox}

\subsection{Stopping a Session}

Press \textbf{Ctrl+C} on the operator laptop.  There is no automatic watchdog in the
firmware --- when packets stop arriving, the ESP32 continues targeting the last
commanded steering angle, throttle, and brake values.  The safety driver must apply
the physical foot brake immediately when the remote session ends.

% ══════════════════════════════════════════════════════════════════════════════
\section{How It Works}
% ══════════════════════════════════════════════════════════════════════════════

\subsection{Step 1 --- Car Registers with the Relay Server}

On boot the ESP32 sends a \texttt{HEARTBEAT} UDP packet to the relay server every
2\,s.  This registers the ESP32's public IP/port and keeps the hotspot NAT mapping
alive so the server can route packets back.

\subsection{Step 2 --- Operator Sends Control Commands}

\texttt{g29\_to\_server\_original.py} polls the G29 at 60\,Hz and sends:
\begin{center}\texttt{CMD;\textit{seq};\textit{steer};\textit{throttle};\textit{brake}}\end{center}

\begin{center}
\begin{tabular}{lll}
  \toprule
  \textbf{Field} & \textbf{Range} & \textbf{Description}\\
  \midrule
  \texttt{steer}    & $-1.0\ldots+1.0$ & $-1.0=$ full left, $+1.0=$ full right\\
  \texttt{throttle} & $0.0\ldots1.0$   & $0.0=$ released, $1.0=$ fully pressed\\
  \texttt{brake}    & $0.0\ldots1.0$   & $0.0=$ released, $1.0=$ fully pressed\\
  \bottomrule
\end{tabular}
\end{center}

\subsection{Step 3 --- Relay Server Forwards to the Car}

\texttt{relay\_server.py} receives \texttt{CMD} packets from the operator laptop and
forwards them verbatim to the last registered \texttt{HEARTBEAT} address.

\subsection{Step 4 --- ESP32 Drives the Actuators}

\subsubsection*{Steering (Closed-Loop PID)}

\[
  \underbrace{\texttt{steer}}_{{-1}\ldots{+1}}
  \;\times\; 250
  \;\longrightarrow\;
  \text{targetSTA}
  \;\xrightarrow{\text{low-pass filter}}\;
  \text{smoothTarget}
  \;\xrightarrow{\text{PID + CAN feedback}}\;
  \text{torque}
\]
\begin{align*}
  \text{DAC\,2, Ch\,C} &= 2.47\,\text{V baseline} + \text{torque}\\
  \text{DAC\,2, Ch\,D} &= 2.47\,\text{V baseline} - \text{torque}
\end{align*}
Actual steering angle is read from CAN ID \texttt{0x2} (little-endian \texttt{int16},
divide by 10 for degrees) and fed back every loop.

\subsubsection*{Throttle (Open-Loop)}

\[
  \texttt{throttle}\times 300 = \text{accelMag}\quad(0\text{--}300)
\]
\begin{align*}
  \text{DAC\,1, Ch\,A} &= 0.37\,\text{V} + \text{accelMag}\\
  \text{DAC\,1, Ch\,B} &= 0.75\,\text{V} + 2\times\text{accelMag}
\end{align*}

\subsubsection*{Brake (Open-Loop)}

\[
  \texttt{brake}\times 900 = \text{brakeMag}\quad(0\text{--}900)
\]
\begin{align*}
  \text{DAC\,1, Ch\,C} &= 3.42\,\text{V} - \text{brakeMag}\\
  \text{DAC\,1, Ch\,D} &= 1.51\,\text{V} + \text{brakeMag}
\end{align*}

% ══════════════════════════════════════════════════════════════════════════════
\section{Troubleshooting}
% ══════════════════════════════════════════════════════════════════════════════

\begin{longtable}{>{\ttfamily\small}p{4.2cm}p{4.5cm}p{5.5cm}}
  \toprule
  \textbf{\rmfamily Symptom} & \textbf{Likely Cause} & \textbf{Fix}\\
  \midrule\endhead\bottomrule\endfoot
  No joystick found.
    & G29 not detected
    & Unplug/replug G29; wait for calibration spin; rerun.\\[5pt]
  CAR CONNECTED never appears
    & ESP32 cannot reach OCI
    & Confirm hotspot is on; confirm UDP 4210 is open in OCI security list and \texttt{iptables}.\\[5pt]
  Car does not respond
    & CMD packets not reaching ESP32
    & Confirm relay server is running and CAR CONNECTED was printed.\\[5pt]
  Steering drives to limit and stays
    & CAN bus not working; no PID feedback
    & Check MCP2515 wiring; verify \textbf{16\,MHz} crystal on PCB.\\[5pt]
  MCP2515 Init Failed (Serial Monitor)
    & SPI wiring or wrong CS pin
    & Check wiring against Section~3 pin table.\\[5pt]
  Throttle/brake has no effect
    & DAC\,1 values wrong
    & Open Serial Monitor; verify \texttt{accelMag}/\texttt{brakeMag} change when pedals are pressed.\\[5pt]
  Steering oscillates or overshoots
    & PID gains need tuning
    & Reduce \texttt{kp}; increase \texttt{kd} to damp.  Retest on stationary vehicle.\\
\end{longtable}

% ══════════════════════════════════════════════════════════════════════════════
\section{Multi-Camera Streaming with OBS}
% ══════════════════════════════════════════════════════════════════════════════

The car-side setup supports USB cameras for a live video feed.
To broadcast multiple camera angles simultaneously to a video call:

\begin{enumerate}
  \item Install \textbf{OBS Studio} (\url{https://obsproject.com}).
  \item Add each USB camera as a \textbf{Video Capture Device} source and arrange
        them in a single OBS scene (e.g., side-by-side or picture-in-picture).
  \item Enable \textbf{OBS Virtual Camera} via the Controls panel
        (\textit{Start Virtual Camera}).
  \item In Zoom (or any conferencing app), select \textbf{OBS Virtual Camera} as the
        camera source.
\end{enumerate}

The remote participant sees the full multi-camera OBS scene without screen sharing.

\end{document}


